\chapter*{Conclusion}\addcontentsline{toc}{chapter}{Conclusion}
One inventory-control problem links the models developed here. Markovian diffusions expose the role of price scale and portfolio covariance; non-Markovian kernels expose memory; finite Markovian lifts recover recursive simulation; scalar compression makes the HJB tractable while discarding generator information. The full residual makes that last loss measurable instead of equating calibration with closure.

The results are therefore read through a common accounting and a hierarchy of approximations, rather than as three independent demonstrations of superiority. Inventory-aware quoting reduces exposure in the diffusion and historical finite-lift samples; the positive baseline has favorable paired liquidated-PnL results in its declared E1/E2 replay. Adding the covariance proxy lowers integrated price risk but lowers sample PnL and certainty equivalent and increases the reconstructed full residual. The bootstrap interval includes zero, so the sample CE decrease is not a population-loss conclusion.

These outcomes motivate a correction method with separate numerical, residual and economic gates. The current contribution is a connected model hierarchy, reproducible baseline and comparative diagnostic, not a fitted P3, a convergence certificate or a latency gain. The remaining independent-reference and source-verification tasks delimit what can be concluded from the present study.
